% !TeX root = closed-systems-from-comparison-completeness.tex
% ============================================================
\chapter*{Orientation to Part VIII}
\phantomsection
\label{part:intro-observable}
% ============================================================
\begin{chapterorientation}
\orientationitem{Settled}{The main closure, transport, geometry, phase, and scalar-sector constructions have been completed.}
\orientationitem{Task}{Part~VIII asks how much state-side observable content is determined by the realized regular sector.}
\orientationitem{Handoff}{The output is a marked extension: raw inexhaustibility, regular-sector identification, and the automorphism bridge problem.}
\end{chapterorientation}

Part~VIII is not a hidden premise for the preceding chapters.  It is a
state-side continuation of the completed stack.  Its purpose is to
separate what is already readable from the regular realized sector from
what remains an explicit bridge problem for dynamics.

\Cref{ch:dynamics-program} proves the raw inexhaustibility result,
identifies the regular sector, and states the automorphism bridge as the
next structural target.  The chapter is therefore positioned as an
extension of the monograph's theorem chain, not as a foundation for the
earlier closure, curvature, phase, or numerical results.
